\documentclass[11pt]{article}
\usepackage{color}

%----------Packages----------
\usepackage{amsmath}
\usepackage{amssymb}
\usepackage{amsthm}
\usepackage{dsfont}
\usepackage{mathrsfs}
\usepackage{stmaryrd}
\usepackage[all]{xy}
\usepackage[mathcal]{eucal}
\usepackage{verbatim}  %%includes comment environment
\usepackage{fullpage}  %%smaller margins
\usepackage{multicol}
\usepackage{hyperref}
\usepackage{graphicx}

%Packages added to get the code block
\usepackage{listings}
\usepackage{xcolor}
\usepackage{tikz-cd}

\definecolor{codegreen}{rgb}{0,0.6,0}
\definecolor{codegray}{rgb}{0.5,0.5,0.5}
\definecolor{codepurple}{rgb}{0.58,0,0.82}
\definecolor{backcolour}{rgb}{0.95,0.95,0.92}

\lstdefinestyle{mystyle}{
    backgroundcolor=\color{backcolour},   
    commentstyle=\color{codegreen},
    keywordstyle=\color{magenta},
    numberstyle=\tiny\color{codegray},
    stringstyle=\color{codepurple},
    basicstyle=\ttfamily\footnotesize,
    breakatwhitespace=false,         
    breaklines=true,                 
    captionpos=b,                    
    keepspaces=true,                 
    numbers=left,                    
    numbersep=5pt,                  
    showspaces=false,                
    showstringspaces=false,
    showtabs=false,                  
    tabsize=2
}

\lstset{style=mystyle}


%---------Commands----------
\def\ZZ{\mathbb{Z}}
\def\QQ{\mathbb{Q}}
\def\PP{\mathbb{P}}
\def\NN{\mathbb{N}}
\def\OO{\mathcal{O}}
\def\AA{\mathbb{A}}
\def\BB{\mathbb{B}}
\def\FF{\mathbb{F}}
\def\LL{\mathbb{L}}
\def\RR{\mathbb{R}}
\def\CC{\mathbb{C}}
\def\GG{\mathbb{G}}
\def\lcm{\mathrm{lcm}}
\def\ord{\mathrm{ord}}
\def\rep{\mathrm{Rep}}
\def\dom{\mathrm{Dom}}
\def\ve{\varepsilon}
\def\gal{\mathrm{Gal}}
\def\emb{\mathrm{Emb}}
\def\inter{\mathrm{Int}}
\def\sub{\mathrm{Sub}}

\newcommand{\dlegendre}[2]{\ensuremath{\left( \frac{#1}{#2} \right) }}
\newcommand{\tlegendre}[2]{\ensuremath{\big( \frac{#1}{#2} \big) }}

%----------Theorems----------
\theoremstyle{definition}
\newtheorem{theorem}{Theorem}[section]
\newtheorem{proposition}[theorem]{Proposition}
\newtheorem{lemma}[theorem]{Lemma}
\newtheorem{corollary}[theorem]{Corollary}
\newtheorem{definition}[theorem]{Definition}
\newtheorem{example}[theorem]{Example}
\newtheorem*{definition*}{Definition}
\newtheorem{nondefinition}[theorem]{Non-Definition}
\newtheorem{problem}{Problem}[section]
\newtheorem*{tip}{Pro Tip}
\newtheorem*{remark}{Remark}
\numberwithin{equation}{subsection}

\newenvironment{solution}
{%
  \vspace{0.5em}%
  \hrule%
  \vspace{0.2em}%
  \noindent\textbf{Solution:}\!\!\!\!\!
  \pushQED{\qed}
}
{%
  \popQED
  \vspace{.2em}
  \hrule%
  \vspace{0.75em}%
}

\newcommand{\separate}{
    \vspace{.2in}
    \hrule
    \vspace{.2in}
}


\newcommand{\head}[1]{
	\begin{center}
		{\large #1}
		\vspace{.05 in}
	\end{center}
	\bigskip 
}

\newcommand{\blue}[1]{\textcolor{blue}{#1}}
\newcommand{\red}[1]{\textcolor{red}{#1}}

%Problem Set Data
%number = 8
%course = 261
%semester = Fall 2025
%path = /Users/thyde/Documents/Teaching/Vassar/Semesters/2025 Fall/Math 261 Fall 2025/Problem Sets/Problem Set 8

\begin{document}
\pagestyle{plain}

%%---  sheet number for theorem counter
\setcounter{section}{0}   

\head{MATH 999, Summer 2026\\ Side Note 0: Quotient by Sums\\
Trevor Hyde}

%HEADER

\begin{center}
   A number will find\\
   fulfillment enough\\
   in knowing its mind\\
   and doing its stuff.\\
   \vspace{.1in}

   \emph{Numbers} by Piet Hein
\end{center}
\vspace{.1in}

In our first meeting we discussed the problem of determining which ideals in $\ZZ[x]$ were maximal.
The trick for testing maximality was to take a candidate ideal $J$, determine the quotient $\ZZ[x]/J$, and see if the ring we get is a field.
For ideals like $J = \langle x^2 + 1, 2\rangle$ I asserted that we could work out the quotient $\ZZ[x]/\langle x^2 +1, 2\rangle$ by taking the quotient in stages,
\[
    \ZZ[x]/\langle x^2 +1, 2\rangle \cong (\ZZ[x]/\langle x^2 + 1\rangle)/\langle 2\rangle
    \cong \ZZ[i]/\langle 2 \rangle .
\]
Since you likely didn't learn this in your algebra class, I wanted to explain the rigorous reasoning behind this and show you some neat examples and consequences.

Note that when we have an ideal like $\langle x^2 + 1, 2\rangle$ generated by more than one element, this is equivalent to expressing our ideal as a sum $\langle x^2 + 1, 2\rangle = \langle x^2 + 1\rangle + \langle 2 \rangle$.
Hence what we need is to analyze quotients like $R/(I + J)$ for a ring $R$ and ideals $I, J \subseteq R$.
Recall that we have reduction modulo $I$ homomorphism $\rho_I : R \to R/I$ defined by $\rho_I(a) := a + I$ or equivalently $\rho_I(a) := a \bmod I$.
Also recall that if $\rho : R \to S$ is a surjective ring homomorphism and if $J \subseteq R$ is an ideal, then $\rho(J) \subseteq S$ is an ideal.
(\blue{If you don't remember this, work through the proof.
It's a quick one, but useful for reminding yourself why we need the surjective hypothesis.})

\begin{proposition}
    Let $R$ be a commutative ring and let $I, J \subseteq R$ be ideals.
    Then
    \[
        R/(I + J) \cong (R/I)/\rho_I(J).
    \]
\end{proposition}

\begin{proof}
    Whenever we want to prove an isomorphism between quotient rings, you should immediately reach for the first isomorphism theorem.
    We start by constructing a surjective homomorphism $\phi : R \to (R/I)/\rho_I(J)$ and then we check that its kernel is precisely $I + J$ and we win.

    In this case, our hands are tied: there is only one natural map $\phi$ to construct.
    First off we have $\rho_I : R \to R/I$ and then we have $\rho_{\rho_I(J)} : R/I \to (R/I)/\rho_I(J)$.
    Both of these maps are surjective, hence so is their composite $\phi : = \rho_{\rho_I(J)}\rho_I$.
    I claim that $\ker\phi = I + J$.
    Let's unpack the definition of $\phi$.
    Suppose $r \in R$, then
    \[
        \phi(r) = \rho_{\rho_I(J)}\rho_I(r)
        = \rho_{\rho_I(J)}(r + I)
        = (r + I) + \rho_I(J).
    \]
    Note that $\rho_I(J) = \{j + I : j \in J\} = I + J$.
    Hence
    \[
        (r + I) + \rho_I(J) = (r + I) + (I + J) = r + I + J.
    \]
    Thus $r \in \ker\phi$ if and only if $r \in I + J$.
    Therefore $\ker\phi = I + J$ and the first isomorphism theorem implies that we have an isomorphism
    \[
        \bar\phi : R/(I + J) \to (R/I)/\rho_I(J).\qedhere
    \]
\end{proof}

The ideal $\rho_I(J)$ is really just the same ideal $J$, but considered within the quotient ring $R/I$.
Hence it is common to abuse notation and simply write $J$ instead of $\rho_I(J)$.
That gives us the simpler looking identity
\[
    R/(I + J) \cong (R/I)/J.
\]
This proposition justifies the calculation we did at the top of this note and during the first meeting.
Using induction this generalizes to sums of more than two ideals.

Since addition of ideals is commutative, we get the following identity
\[
    (R/I)/J \cong (R/J)/I.
\]
This doesn't look like much and formally it is a triviality, but in practice this can be quite profound.

Let's (re)consider\footnote{If you took number theory with me, we already studied this question.} the problem of which primes $p \in \ZZ$ are no longer prime in the Gaussian integers $\ZZ[i]$.
For example, $3$ does not factor in $\ZZ[i]$, hence remains prime, whereas $5 = (2 + i)(2 - i)$ factors in $\ZZ[i]$.
We can detect these factorizations via the structure of the quotient ring $\ZZ[i]/\langle p\rangle$.
If $p$ factors as, say, $p = \alpha \beta$ in $\ZZ[i]$ where $\alpha$ and $\beta$ are not units, then the quotient ring $\ZZ[i]/\langle p\rangle$ will have zero divisors: $\alpha, \beta \not\equiv 0 \bmod p$ but $\alpha\beta \equiv p \equiv 0 \bmod p$.
On the other hand, if $p$ does not factor, then it is irreducible, and $\ZZ[i]/\langle p \rangle$ will be a domain.

Here is the fun trick: We use the isomorphism $\ZZ[i] \cong \ZZ[x]/\langle x^2 + 1\rangle$ to get
\[
    \ZZ[i]/\langle p \rangle \cong (\ZZ[x]/\langle x^2 + 1\rangle)/\langle p \rangle
    \cong (\ZZ[x]/\langle p\rangle)/\langle x^2 + 1\rangle.
\]
Now what is the quotient $\ZZ[x]/\langle p\rangle$?
Since there are no algebraic relations between $x$ and the integers, this quotient is simply
\[
    \ZZ[x]/\langle p\rangle \cong (\ZZ/\langle p\rangle)[x] \cong \FF_p[x].
\]
(\blue{Feel suspicious? Work it out using the definitions.})
Substituting back into our isomorphisms above we get
\[
    \ZZ[i]/\langle p \rangle \cong \FF_p[x]/\langle x^2 + 1\rangle.
\]
Isomorphisms preserve properties like being a domain or being a field.
Thus this isomorphism tells us that \emph{$p$ is irreducible in $\ZZ[i]$ if and only if $x^2 + 1$ is irreducible in $\FF_p[x]$}.
The former question feels like a tough problem about the arithmetic of this exotic ring $\ZZ[i]$, but the latter question is simply a problem of factoring the polynomial $x^2 + 1$ modulo $p$.
Since $x^2 + 1$ is quadratic, it factors in $\FF_p[x]$ if and only if it has a root in $\FF_p$.
That, in turn, is equivalent to saying that $-1$ is a square modulo $p$, which we can express with the Legendre symbol determining $\tlegendre{-1}{p}$.
Namely, $x^2 + 1$ will be irreducible if and only if $\tlegendre{-1}{p} = -1$, and quadratic reciprocity tells us that this is true if and only if $p \equiv -1 \bmod 4$.
Thus we conclude that $p$ is irreducible in $\ZZ[i]$ if and only if $p \equiv -1 \bmod 4$.
Very nice!

This same trick works for all sorts of extensions.
For example, I encourage you to try to figure out which primes $p \in \ZZ$ factor in $\ZZ[\sqrt{5}]$.
\end{document}
